\section{What restimulation transcriptomes say that imaging would need to detect}
\label{sec:partC}
No public dataset images trained or tolerised monocytes, but three public datasets measure
their transcripts after restimulation. We use them to ask which differences an imaging
study would have to reproduce, with three safeguards: every restimulation label is checked
against the GEO sample records and against a positive control (restimulated untrained
cells must exceed unstimulated untrained cells for TNF, IL6 and IL1B in every donor), every
expected direction is quoted from the dataset's source publication, and sample identity is
verified before any observation is reported.

\paragraph{Sourced directions.} Of the \val{pc.v2.n_sourced} TNF and IL-6 directions the
source publications report after restimulation, \val{pc.v2.n_match} agree in sign with the
mean donor-paired difference in the deposited transcripts: dimethyl itaconate (GSE212282) and IL-4 (GSE185433)
increase TNF and IL6 transcripts, as their source papers report for the secreted proteins
\citep{ferreira2023,schrijver2023}. The two exceptions are both $\beta$-glucan.

\paragraph{$\beta$-glucan lowers the TNF transcript at 4\,h after restimulation, in two
independent datasets.} In GSE212282, $\beta$-glucan-trained macrophages restimulated with LPS
for 4\,h on day 6 have less TNF mRNA than untrained cells in all three donors (log2
difference \val{pc.212.tnf.d7.min} to \val{pc.212.tnf.d7.max}), although the source paper
reports more secreted TNF-$\alpha$ after 24\,h of restimulation
\citep{ferreira2023}. Before reporting this we checked sample identity. All
\val{qc.212.n} samples agree on sex within each donor; the three donors carry distinct
presence/absence genotypes at HLA-DRB5 and GSTM1, consistent across every sample of a donor,
so no sample belongs to another donor; and after removing each donor's mean, every
$\beta$-glucan sample's profile is most similar to the other donors' $\beta$-glucan samples
(day-6 restimulated arms: \val{qc.212.d7.match} samples matched to their own arm, the one
exception being a $\beta$-glucan + DMI sample matched to DMI). Over the protocol the
$\beta$-glucan effect on TNF reverses: it raises the transcript during the first exposure
(4\,h: \val{pc.212.tnf.4H.min} to \val{pc.212.tnf.4H.max}; day 1: \val{pc.212.tnf.d1.min} to
\val{pc.212.tnf.d1.max}), is inconsistent at rest on day 6 (\val{pc.212.tnf.d6.min} to
\val{pc.212.tnf.d6.max}) and lowers it after restimulation (Figure~\ref{fig:restim}).

The observation replicates in GSE141656, an independent dataset none of whose three donors
matches a GSE212282 donor at all four genotype-like loci we compared: after 4\,h of heat-killed \emph{M. tuberculosis}
on day 6, $\beta$-glucan-trained macrophages have less TNF mRNA in every donor
(\val{pc.141.TNF.min} to \val{pc.141.TNF.max}; restimulation positive control at least
\val{pc.141.pc.tnf.min} log2 units in every donor), and IL6 (\val{pc.141.IL6.mean}) and
IL1B (\val{pc.141.IL1B.mean}) move the same way. Its source publication, too, reports more
TNF-$\alpha$ protein after 24\,h of restimulation \citep{moorlag2020}. A third dataset with
only two replicates (GSE242947, \citealp{cheng2024dectin}) is consistent at 1\,h of LPS
(\val{pc.242.tnf.1h}) and inconsistent at 3\,h (\val{pc.242.tnf.3h}).

We therefore report a replicated, transcript-level observation: at 4\,h after
restimulation, $\beta$-glucan-trained human macrophages have lower TNF mRNA than untrained
ones, in two datasets and six donors, while the same studies report higher TNF-$\alpha$
protein at 24\,h. One explanation is kinetic, that trained cells reach their TNF transcript
peak earlier and resolve it sooner; it is a hypothesis, and the test is a restimulation time
course from 0.5 to 24\,h with paired RNA and protein. With three donors per dataset no
confidence interval is reported; the smallest achievable exact $p$ for a three-donor
sign-flip test is 0.25, so the evidence is the consistency across six of six donors and two
restimuli, not a $p$-value.
